\chapter{Stress--Strain Relation}

\section*{Introduction}

When you sit on a chair, the seat sags slightly. The amount of sagging depends on the material's stiffness---a steel chair barely deforms, while a wooden chair flexes more. The stress-strain relation quantifies this: it's the material's ``personality,'' telling you how it responds to being pushed, pulled, or twisted. Engineers select materials based on this relation: steel for bridges (high stiffness), rubber for tires (low stiffness, high elasticity).

\begin{enumerate}
\item \textbf{Hooke's Law (1D):} Stress is proportional to strain in the linear elastic regime.

\item \textbf{Generalized Hooke's Law (3D):} The full tensor relation between stress and strain for an isotropic material.

\end{enumerate}


\[
\sigma = E\varepsilon
\]

\section*{Fundamental Equations Used}

The derivation starts from Hooke's law for a single spring, $F = -k\Delta L$, generalised to a continuum.

\section*{Assumptions}

\textbf{Assumptions:} The material is linear elastic (stress proportional to strain). Deformations are small. The material is homogeneous and isotropic.


\textbf{Limitations:} Fails beyond the yield point (plastic deformation). Does not describe viscoelastic materials (polymers), which have time-dependent strain. Anisotropic materials (composites, wood) require tensor forms.


\section*{Mathematical Derivation}

\textbf{Step 1: Hooke's Law for a Single Spring.}

Consider a spring with rest length $L_0$. When stretched by an amount $\Delta L$, it exerts a restoring force $F = -k\,\Delta L$, where $k$ is the spring constant. The negative sign indicates the force opposes the displacement. This is the simplest statement of linear elasticity: the restoring force is proportional to the deformation.

\textbf{Step 2: From Force to Stress and Strain.}

Now consider a uniform bar of cross-sectional area $A$ and original length $L_0$. Apply a tensile force $F$ to one end. We define two quantities that remove the dependence on geometry:

\begin{align}
\text{Stress:} \quad \sigma &= \frac{F}{A}, \\
\text{Strain:} \quad \varepsilon &= \frac{\Delta L}{L_0}.
\end{align}

Stress measures the intensity of internal forces (force per unit area); strain measures the fractional deformation (dimensionless). Both are intensive quantities that characterize the material, not the specimen.

\textbf{Step 3: The 1D Stress--Strain Relation.}

For small deformations, experiments show that stress is proportional to strain:

\begin{equation}
\boxed{\sigma = E\varepsilon}
\end{equation}

where $E$ is Young's modulus (the modulus of elasticity). This is the macroscopic manifestation of Hooke's law: $E = kL_0/A$ connects the spring constant to the material property. The linear regime ends at the yield stress $\sigma_Y$, beyond which plastic (permanent) deformation occurs.

\textbf{Step 4: Lateral Contraction---Poisson's Ratio.}

When a bar is stretched axially, it also contracts laterally. If the axial strain is $\varepsilon_{xx}$, the transverse strains are:

\begin{equation}
\varepsilon_{yy} = \varepsilon_{zz} = -\nu\,\varepsilon_{xx},
\end{equation}

where $\nu$ is Poisson's ratio. For most metals, $\nu \approx 0.3$. For a perfectly incompressible material, $\nu = 0.5$. The physical constraint of positive strain energy requires $-1 < \nu < 0.5$.

\textbf{Step 5: Generalized Hooke's Law in 3D.}

For a general 3D stress state, we write the strain components as linear functions of all stress components. For an isotropic material, symmetry and the requirements of objectivity reduce the number of independent elastic constants to two. Using Lam\'e parameters $\lambda$ and $\mu$:

\begin{align}
\varepsilon_{xx} &= \frac{1}{E}\left[\sigma_{xx} - \nu(\sigma_{yy} + \sigma_{zz})\right], \\
\varepsilon_{yy} &= \frac{1}{E}\left[\sigma_{yy} - \nu(\sigma_{xx} + \sigma_{zz})\right], \\
\varepsilon_{zz} &= \frac{1}{E}\left[\sigma_{zz} - \nu(\sigma_{xx} + \sigma_{yy})\right], \\
\varepsilon_{xy} &= \frac{1 + \nu}{E}\,\sigma_{xy} = \frac{\sigma_{xy}}{2\mu}, \\
\varepsilon_{yz} &= \frac{1 + \nu}{E}\,\sigma_{yz} = \frac{\sigma_{yz}}{2\mu}, \\
\varepsilon_{zx} &= \frac{1 + \nu}{E}\,\sigma_{zx} = \frac{\sigma_{zx}}{2\mu}.
\end{align}

\textbf{Step 6: Compact Tensor Notation.}

In Voigt notation, or equivalently in direct tensor notation, the generalized Hooke's law is:

\begin{equation}
\varepsilon_{ij} = \frac{1+\nu}{E}\,\sigma_{ij} - \frac{\nu}{E}\,\sigma_{kk}\,\delta_{ij},
\end{equation}

where $\sigma_{kk} = \sigma_{xx} + \sigma_{yy} + \sigma_{zz}$ is the trace (the sum of normal stresses, equal to $-3p$ for hydrostatic pressure) and $\delta_{ij}$ is the Kronecker delta. The inverse form, expressing stress in terms of strain, uses Lam\'e parameters:

\begin{equation}
\sigma_{ij} = 2\mu\,\varepsilon_{ij} + \lambda\,\varepsilon_{kk}\,\delta_{ij},
\end{equation}

where $\lambda = \frac{E\nu}{(1+\nu)(1-2\nu)}$ and $\mu = \frac{E}{2(1+\nu)}$ is the shear modulus.

\textbf{Step 7: Physical Meaning of the Two Constants.}

The two independent constants ($E$ and $\nu$, or equivalently $\lambda$ and $\mu$) encode all elastic behavior of an isotropic material:

\begin{itemize}
\item The bulk modulus $K = E/[3(1-2\nu)] = \lambda + \frac{2}{3}\mu$ measures resistance to uniform compression.
\item The shear modulus $\mu = E/[2(1+\nu)]$ measures resistance to shape change at constant volume.
\item The Young's modulus $E = 9K\mu/(3K + \mu)$ measures resistance to uniaxial stretching.
\end{itemize}

Any two of these three moduli ($K$, $\mu$, $E$) determine the others, confirming that two independent constants suffice for an isotropic linear elastic material.

\section*{Characteristics of the Equation}

The two independent elastic constants can be expressed as: Young's modulus $E$, Poisson's ratio $\nu$, shear modulus $\mu = E/2(1+\nu)$, bulk modulus $K = E/3(1-2\nu)$. For stable isotropic materials: $-1 < \nu < 0.5$ and $E > 0$.
\section*{Solution Methods}

For simple loading (uniaxial tension, pure shear): direct application. For complex geometries: combine with stress equilibrium and compatibility equations, or use FEA. Experimental determination: tensile tests, ultrasonic measurements, nanoindentation.
\section*{Verifications}

Tensile testing machines measure stress-strain curves directly. Ultrasound measures wave speeds, from which elastic moduli can be calculated. Bending tests on beams give $E$ from deflection measurements.
\section*{Applications}

Material selection for engineering design, structural analysis, quality control (measuring $E$ and $\nu$), seismic wave speed calculation ($v_p$ and $v_s$ depend on $\lambda$ and $\mu$), biomedical implants (matching bone stiffness to avoid stress shielding).
\section*{What Richard Feynman Would Have Asked}

\subsection*{1. The Plain-English Translation}
\textit{``What does it really mean for a material to have a Young's modulus?''}

The Core Meaning: Young's modulus tells you how much force is needed to stretch a material to a given fraction of its original length. A material with $E = 200$ GPa (steel) requires twice the force to stretch by the same percentage as one with $E = 100$ GPa. It is the material's ``resistance to being pulled apart,'' expressed as a ratio of stress to strain.

The Metaphor: Think of materials as springs with different stiffness. A steel spring is stiff (high $E$): you need to push hard to compress it a little. A rubber spring is soft (low $E$): a small force produces a large deformation. The stress-strain relation is just Hooke's law for the material itself.

\subsection*{2. The Localized Mechanics}
\textit{``Look at the microstructure of the material. What determines $E$?''}

He would press you on the atomic origin of stiffness. Young's modulus is determined by the curvature of the interatomic potential energy at the equilibrium spacing. Stronger atomic bonds (covalent, ionic) give higher $E$; weaker bonds (van der Waals) give lower $E$. Diamond ($E \approx 1000$ GPa) has strong covalent bonds; rubber ($E \approx 0.01$ GPa) has weak cross-links between polymer chains.

The Smoothness Check: He would ask whether $E$ is truly constant. At very high strains, the linear approximation breaks down---the stress-strain curve curves upward (strain hardening) or downward (necking). Even in the ``linear'' regime, $E$ varies slightly with temperature and strain rate.

\subsection*{3. Testing the Extreme Limits}
\textit{``What happens at extreme strains or temperatures?''}

The Yield Point: Beyond $\sigma_{\text{yield}}$, permanent (plastic) deformation occurs. The stress-strain relation is no longer linear---the material has a ``memory'' of its loading history. This is why you can bend a paperclip back and forth until it breaks: each cycle accumulates plastic strain.

The Melting Point: At the melting temperature, the crystalline structure collapses and $E \to 0$. The material flows like a liquid---it has no resistance to shear (only to compression). This is why volcanologists model magma as a fluid, not a solid.

The Quantum Limit: At absolute zero, thermal vibrations cease but quantum zero-point motion remains. The elastic constants approach finite values determined entirely by the electronic structure.

\subsection*{4. Dimensionality and Geometry}
\textit{``Does a thin wire have the same $E$ as a thick block of the same material?''}

He would point out that $E$ is an intrinsic material property---it doesn't depend on geometry. But the apparent stiffness of a structure (how much it deflects under a given load) depends strongly on geometry. A thin wire stretches more than a thick block under the same force, even though $E$ is the same. The stress-strain relation separates material behavior from geometric effects.

\subsection*{5. Cross-Disciplinary Connection}
\textit{``Where else does this same linear proportionality appear?''}

The same stress-strain structure appears in: Ohm's law (current proportional to voltage, $I = V/R$), Hooke's law for springs ($F = kx$), Newton's law of viscosity ($\tau = \mu \dot{\gamma}$), and even in economics (supply/demand curves near equilibrium). The mathematical structure---a linear relation between a ``response'' and a ``driving force''---is the hallmark of linear physics near equilibrium.
\section*{FAQ}

\textbf{Q: What happens if $\nu > 0.5$?}\\
\textbf{A:} The bulk modulus becomes negative, meaning the material would expand under hydrostatic compression---physically impossible for stable materials. Rubber has $\nu \approx 0.499$ (nearly incompressible), but never exceeds 0.5.

\textbf{Q: Is steel or aluminum stiffer?}\\
\textbf{A:} Steel ($E \approx 200$ GPa) is about three times stiffer than aluminum ($E \approx 70$ GPa). But stiffness isn't everything---aluminum is lighter, so for weight-critical applications, the strength-to-weight ratio matters more.
\section*{Important Bibliography}

\begin{enumerate}
\item Timoshenko, S.P. and Goodier, J.N., \textit{Theory of Elasticity}, 3rd ed. (McGraw-Hill, 1970), Chapters 1--3.
\item Landau, L.D. and Lifshitz, E.M., \textit{Theory of Elasticity}, Vol. 7 of \textit{Course of Theoretical Physics}, 3rd ed. (Butterworth-Heinemann, 1986), Chapters 1--2.
\item Hooke, R., \textit{De Potentia Restitutiva} (John Martyn, 1678).
\item Love, A.E.H., \textit{A Treatise on the Mathematical Theory of Elasticity}, 4th ed. (Cambridge University Press, 1927), Chapter 1.
\item Johnson, K.L., \textit{Contact Mechanics} (Cambridge University Press, 1985), Chapter 2.
\end{enumerate}

